\documentclass[10pt,a4paper]{amsart}
\usepackage[margin=19mm]{geometry}
\usepackage{amsmath,amssymb,mathtools}
\usepackage[scr=rsfs,cal=euler]{mathalfa}
\usepackage{xfrac}
\usepackage[unicode,hidelinks,bookmarks=false]{hyperref}
\newcommand{\ie}{i.e.\ }
\newcommand{\cf}{cf.\ }
\newcommand{\resp}{respectively\ }
\newcommand{\Subsection}{\subsection}
\providecommand{\nobreakdash}{\nobreak\hbox{-}}
\setlength{\emergencystretch}{2em}
\multlinegap0pt
\usepackage{amssymb}
\usepackage{subfig}
\usepackage{xcolor}
\usepackage{mathtools}
\usepackage{enumerate}
\newtheorem{lemma}{Lemma}
\newcommand{\bsk}{\boldsymbol{k}}
\newcommand{\bsone}{\boldsymbol{1}}
\newcommand{\dunif}{\mathbb{U}}
\newcommand{\e}{\mathbb{E}}
\renewcommand{\geq}{\geqslant}
\renewcommand{\leq}{\leqslant}
\newcommand{\natu}{\mathbb{N}}
\newcommand{\rd}{\,\mathrm{d}}
\newcommand{\var}{\mathrm{Var}}
\title{Quasi-Monte Carlo confidence intervals using quantiles of randomized nets}
\author{Zexin Pan}
\date{Source: arXiv:2504.19138v2 (CC BY 4.0)}
\begin{document}
\maketitle

\noindent\emph{Source and licence.} Quasi-Monte Carlo confidence intervals using quantiles of randomized nets. Version arXiv:2504.19138v2 (CC BY 4.0), distributed by arXiv under the Creative Commons Attribution 4.0 International licence, http://creativecommons.org/licenses/by/4.0/, as recorded in the arXiv licence metadata for this exact version and shown on the abstract page https://arxiv.org/abs/2504.19138v2. Full licence notice: \texttt{licenses/arxiv-2504.19138-CC-BY-4.0.txt}.\par\smallskip

\noindent\textit{Editorial scope.} Counted unit: the article's lemma lem:sumk (source0.tex lines 365--370). The article's complete original proof of it is delivered with it (source0.tex lines 1131--1214, one \texttt{proof} environment of 84 lines, which the article places away from the statement). No mathematics is written, repaired or re-derived: the statement and the proof are the article's own lines, in the article's own notation, and no retained line is substituted or re-flowed.  Retained uncounted as the source-local context the proof needs: lines 1055--1059; lines 1104--1106; lines 1127--1129.  Nothing was omitted from any retained span, so the package carries no omission record.  No retained line contains a citation, so the wrapper carries no bibliography and no bibliography entry.  No retained line contains a figure environment, an image-inclusion command or drawing code, so no image is delivered and the package declares no asset dependency.  Editorial formatting only, in addition: an AMS \texttt{amsart} preamble that restates the article's own declarations and macros used by the retained lines, namely the environments \texttt{lemma} on separate counters, together with the standard packages those lines need.  The retained environments print in this document as Lemma 1 in order of appearance, whereas the article prints the same items with the numbering of its own full document.  The complete native source of the article as distributed in the arXiv e-print for this exact version is bundled unchanged as \texttt{sources/177149/source0.tex}.
\begin{lemma}\label{lem:likehoodratio} 
When $N\geq 1$,
    $$\frac{1}{|Q_N|q_N^N}\prod_{\ell=1}^\infty (1+q_N^\ell)^s\leq A_s N^{1/4}$$
    with $A_s$ a constant depending on $s$.
\end{lemma}

    \begin{equation}\label{eqn:varLasym}
        \var^L(|\kappa_j|)\sim \frac{1}{\pi}\sqrt{\frac{3 N}{s}}+O(1).
    \end{equation}

\begin{equation}\label{eqn:rversionPrlePrL}
    \Pr((\bsk_1,\dots,\bsk_r)\in A)\leq {\textstyle\Pr}^L((\bsk_1,\dots,\bsk_r)\in A)\Bigl(\frac{1}{|Q_N|q_N^N}\prod_{\ell=1}^\infty (1+q_N^\ell)^s\Bigr)^r.
\end{equation}

\begin{lemma}\label{lem:sumk}
Let $N\geq 1$ and $\bsk_1,\dots,\bsk_r$ be sampled independently from $\dunif (Q_N)$. Then there exist positive constants $A_s, B_s$ depending on $s$ such that for all $r\geq 2$
\begin{equation*}
    \Pr\left(\oplus_{i=1}^r\bsk_i\in Q_N\right)\leq A^r_s N^{r/4} r^{-B_s\sqrt{N}}.
\end{equation*}
\end{lemma}

\begin{proof}[Proof of Lemma \ref{lem:sumk}]
To simplify our notation, we write $\bsk^\oplus=\oplus_{i=1}^r\bsk_i$ with components $(k^\oplus_1,\dots, k^\oplus_s)$. By equation~\eqref{eqn:rversionPrlePrL} and Lemma~\ref{lem:likehoodratio},
    \begin{equation}\label{eqn:PrlePrLforsumk}
        \Pr\big(\bsk^\oplus\in Q_N\big)\leq A^r_s N^{r/4}{\textstyle\Pr}^L\big(\bsk^\oplus\in Q_N\big)\leq A^r_s N^{r/4}{\textstyle\Pr}^L\big(\Vert\bsk^\oplus\Vert_1\leq N\big).
    \end{equation}
    By the definition of $\bsk^\oplus$, $X^\oplus_{j\ell}=\bsone\{\ell\in \kappa^\oplus_j\}$ equals $1$ if and only if $\ell\in \kappa_j$ for an odd number of $\bsk$ among $\bsk_1,...,\bsk_r$.
    By a binomial distribution with success probability $q_N^\ell/(1+q_N^\ell)$,
    \begin{align}\label{eqn:oddbinomial}
        {\textstyle\Pr}^L(X^\oplus_{j\ell}=1)&=\sum_{j=1}^{\lceil r/2\rceil}{r\choose 2j-1}\Big(\frac{q_N^\ell}{1+q_N^\ell}\Big)^{2j-1}\Big(\frac{1}{1+q_N^\ell}\Big)^{r-2j+1}\\
        &=\frac{1}{2}-\frac{1}{2}\Bigl(\frac{1-q_N^\ell}{1+q_N^\ell}\Bigr)^r. \nonumber
    \end{align}
    Also notice that $\{X^\oplus_{j\ell}, j\in1{:}s, \ell\in \natu\}$ are jointly independent under $L(\bsk)$ and $$\Vert\bsk^\oplus\Vert_1=\sum_{j=1}^s\sum_{\ell\in\natu}\ell X^\oplus_{j\ell}.$$
 By Markov's inequality, for any $t>0$
    \begin{align}\label{eqn:chernoffbound}
        {\textstyle\Pr}^L\Big(\Vert\bsk^\oplus\Vert_1\leq N\big)&={\textstyle\Pr}^L\big(\exp\Big(-t\sum_{j=1}^s\sum_{\ell\in\natu}\ell X^\oplus_{j\ell}\Big)\geq e^{-tN}\Big)\\
        & \leq e^{tN} \e^L\Big[\exp\Big(-t\sum_{j=1}^s\sum_{\ell\in\natu}\ell X^\oplus_{j\ell}\Big)\Big]\nonumber\\
        & =e^{tN}\prod_{j=1}^s\prod_{\ell\in N}\Bigl(1-{\textstyle\Pr}^L(X^\oplus_{j\ell}=1)(1-e^{-t\ell })\Bigr)\nonumber\\
        &\leq \exp\Bigl(tN-s\sum_{\ell\in\natu}{\textstyle\Pr}^L(X^\oplus_{j\ell}=1)(1-e^{-t\ell })\Bigr).\nonumber
    \end{align}
    Because ${\textstyle\Pr}^L(X^\oplus_{j\ell}=1)$ is monotonically increasing in $r$ and $r\geq 2$,
    \begin{align*}
        \sum_{\ell\in\natu}{\textstyle\Pr}^L(X^\oplus_{j\ell}=1)(1-e^{-t\ell })
        &\geq \sum_{\ell\in\natu}\Bigl(\frac{1}{2}-\frac{1}{2}\Bigl(\frac{1-q_N^\ell}{1+q_N^\ell}\Bigr)^2\Bigr)(1-e^{-t\ell })\\
        &=\sum_{\ell\in\natu}\frac{1}{2}(1-e^{-t\ell })\frac{(1+q_N^\ell)^2-(1-q_N^\ell)^2}{(1+q_N^\ell)^2}\\
        &=\sum_{\ell\in\natu}2(1-e^{-t\ell })\frac{q_N^\ell}{(1+q_N^\ell)^2}.
    \end{align*}
    Setting $t=-\alpha\log(q_N)$ for $\alpha> 0$, which we will tune later, we have
    \begin{align*}
    \sum_{\ell\in\natu}2(1-e^{-t\ell })\frac{q_N^\ell}{(1+q_N^\ell)^2}
        &=2\sum_{\ell\in\natu}\frac{q_N^\ell}{(1+q_N^\ell)^2}-2\sum_{\ell\in\natu}\frac{q_N^{\alpha\ell}q_N^\ell}{(1+q_N^\ell)^2}.
    \end{align*}
    Similar to equation~\eqref{eqn:varLasym}, because both $q_N^\ell/(1+q_N^\ell)^2$ and $q_N^{\alpha\ell}q_N^\ell/(1+q_N^\ell)^2$ are monotonically decreasing in $\ell$ over $\ell\geq 0$,
    \begin{align*}
       \sum_{\ell\in\natu}\frac{q_N^\ell}{(1+q_N^\ell)^2}\sim &\int_0^\infty \frac{ \exp(-\pi\ell\sqrt{s/12N})}{(1+\exp(-\pi\ell\sqrt{s/12N}))^2}\rd\ell+O(1)\\
       =&\frac{1}{\pi}\sqrt{\frac{12 N}{s}}\int_0^\infty \frac{ \exp(-\ell)}{(1+\exp(-\ell))^2}\rd\ell+O(1)
    \end{align*}
    and
    \begin{align*}
        \sum_{\ell\in\natu}\frac{q_N^{\alpha\ell}q_N^\ell}{(1+q_N^\ell)^2}\sim &\int_0^\infty \frac{ \exp(-\pi(\alpha+1)\ell\sqrt{s/12N})}{(1+\exp(-\pi\ell\sqrt{s/12N}))^2}\rd\ell+O(1)\\
        =&\frac{1}{\pi}\sqrt{\frac{12 N}{s}}\int_0^\infty \frac{ \exp(-(\alpha+1)\ell)}{(1+\exp(-\ell))^2}\rd\ell+O(1). 
    \end{align*}
    Combining $t=-\alpha\log(q_N)=\alpha \pi\sqrt{s/12N}$ with the above equations, we get
    $$tN-s\sum_{\ell\in\natu}{\textstyle\Pr}^L(X^\oplus_{j\ell}=1)(1-e^{-t\ell })\leq tN-2s\sum_{\ell\in\natu}\frac{q_N^\ell-q_N^{\alpha\ell}q_N^\ell}{(1+q_N^\ell)^2}\sim c(\alpha)\sqrt{sN}+O(s)$$
    if $c(\alpha)\neq 0$ with 
    $$c(\alpha)=\alpha\frac{\pi}{\sqrt{12}}-\frac{4\sqrt{3 }}{\pi}\int_0^\infty \frac{\exp(-\ell)- \exp(-(\alpha+1)\ell)}{(1+\exp(-\ell))^2}\rd\ell.$$
    Because $c(\alpha)\to\infty$ as $\alpha\to \infty$ and
    $$c'(\alpha)=\frac{\pi}{\sqrt{12}}-\frac{4\sqrt{3}}{\pi}\int_0^\infty \frac{ \ell\exp(-(\alpha+1)\ell)}{(1+\exp(-\ell))^2}\rd\ell$$
    is strictly increasing in $\alpha$, we see $c(\alpha)$ has a unique minimum $\alpha^*$ over $\alpha\geq 0$. Furthermore,
     $$c'(0)=\frac{\pi}{\sqrt{12}}-\frac{4\sqrt{3}}{\pi}\int_0^\infty \frac{ \ell\exp(-\ell)}{(1+\exp(-\ell))^2}\rd\ell=\frac{\pi}{\sqrt{12}}-\frac{4\sqrt{3}\log(2)}{\pi}<0,$$
     so $\alpha^*>0$ and $c(\alpha^*)<0$. A numerical approximation using Mathematica shows $\alpha^*\approx 0.24$ and $c(\alpha^*)< -0.066$.
By choosing $t=-\alpha^*\log(q)$, we have shown
    $$\exp\Bigl(tN-s\sum_{\ell\in\natu}{\textstyle\Pr}^L(X^\oplus_{j\ell}=1)(1-e^{-t\ell })\Bigr)\leq \exp\Big(c(\alpha^*)\sqrt{sN}+O(s)\Big).$$
    Putting together equation~\eqref{eqn:PrlePrLforsumk} and equation~\eqref{eqn:chernoffbound}, we get
    $$\Pr\big(\bsk^\oplus\in Q_N\big)\leq A^r_sN^{r/4}\exp\Big(c(\alpha^*)\sqrt{sN}+O(s)\Big).$$
    For a threshold $R_s\geq 2$ that we will determine later, we can choose $B_s$ small enough so that $B_s \log(R_s)\leq -c(\alpha^*)\sqrt{s}$. By increasing $A_s$ if necessary to account for the $\exp(O(s))$ term, the conclusion holds for all $N\geq 1$ and $r\leq R_s$.

        It remains to show the conclusion holds for some $A_s, B_s>0$ when $r> R_s$ for some threshold $R_s$. Let $\ell^*$ be the largest $\ell\in \natu$ for which ${\textstyle\Pr}^L(X^\oplus_{j\ell}=1)\geq 1/4$. We conventionally set $\ell^*=0$ if ${\textstyle\Pr}^L(X^\oplus_{j\ell}=1)< 1/4$ for all $\ell\in\natu$. By equation~\eqref{eqn:oddbinomial},
    $${\textstyle\Pr}^L(X^\oplus_{j\ell}=1)=\frac{1}{2}-\frac{1}{2}\Bigl(\frac{1-q_N^\ell}{1+q_N^\ell}\Bigr)^r=\frac{1}{2}-\frac{1}{2}\Bigl(\frac{1-\exp(-\pi\ell\sqrt{s/12N})}{1+\exp(-\pi\ell\sqrt{s/12N})}\Bigr)^r.$$
    Because ${\textstyle\Pr}^L(X^\oplus_{j\ell}=1)$ is monotonically decreasing in $\ell$ over $\ell\geq 0$, $\ell^*$ equals the floor of the solution of ${\textstyle\Pr}^L(X^\oplus_{j\ell}=1)=1/4$. A straightforward calculation gives
    $$\ell^*=\lfloor\log\Bigl(\frac{1+2^{-1/r}}{1-2^{-1/r}}\Bigr)\sqrt{\frac{12N}{\pi^2 s}}\rfloor.$$
    By convexity of the function $f(x)=x^{-1/r}$,
    $$2^{-1-1/r}r^{-1}\leq 1-2^{-1/r}\leq r^{-1}.$$
    Hence
    $$r\leq\frac{1+2^{-1/r}}{1-2^{-1/r}}\leq (1+2^{-1/r})2^{1+1/r}r\leq 8r$$
    and
    \begin{equation}\label{eqn:lstarbound}
      \log(r)\sqrt{\frac{12N}{\pi^2 s}}-1\leq \ell^*\leq \log(8r)\sqrt{\frac{12N}{\pi^2 s}}.  
    \end{equation}
    Using the inequality $1-\exp(-x)\geq x/(1+x)$ when $x\geq 0$, equation~\eqref{eqn:chernoffbound} becomes
    \begin{align}\label{eqn:chernoffbound2}
        {\textstyle\Pr}^L\big(\Vert\bsk^\oplus\Vert_1\leq N\big)&\leq \exp\Bigl(tN-s\sum_{\ell\in\natu}{\textstyle\Pr}^L(X^\oplus_{j\ell}=1)\frac{t\ell}{1+t\ell}\Bigr)\\
    &\leq \exp\Bigl(tN-s\sum_{\ell\in \natu}\bsone\{\ell\leq\ell^*\}\frac{t\ell}{4(1+t\ell)}\Bigr)\nonumber\\
    &= \exp\Bigl(tN-\frac{st\ell^*(\ell^*+1)}{8(1+t\ell^*)}\Bigr).\nonumber
    \end{align}
    Setting $t=\sqrt{s/N}$, we derive from equation~\eqref{eqn:lstarbound} that there exists a large enough $R_s$ so that for all $r\geq R_s$,
    $$1+t\ell^*\leq 1+\log(8r)\sqrt{\frac{12}{\pi^2 }}<2\log(r)\sqrt{\frac{12}{\pi^2 }}$$
    and
    $$st\ell^*(\ell^*+1)\geq  \sqrt{sN}\log(r)\sqrt{\frac{12}{\pi^2 }}\Bigl(\log(r)\sqrt{\frac{12}{\pi^2 }}-\sqrt{\frac{s}{N}}\Bigr)> \sqrt{sN}\log(r)^2\frac{6}{\pi^2}.$$
    By increasing $R_s$ if necessary, we further have for all $r\geq R_s$
    $$tN-\frac{st\ell^*(\ell^*+1)}{8(1+t\ell^*)}\leq \sqrt{sN}-\sqrt{sN}\log(r)\frac{\sqrt{3}}{16\pi}<-\sqrt{sN}\log(r)\frac{\sqrt{3}}{32\pi}.$$
    Putting together equation~\eqref{eqn:PrlePrLforsumk} and equation~\eqref{eqn:chernoffbound2}, we get for $r\geq R_s$
    $$\Pr\big(\bsk^\oplus\in Q_N\big)\leq A^r_sN^{r/4}\exp(-\sqrt{sN}\log(r)\frac{\sqrt{3}}{32\pi}),$$
    which completes the proof.
\end{proof}

\end{document}
